% GENERATED FILE. Edit canonical lessons, then run npm run generate:books.
% canonical-source-hash: fnv1a64:81d3f2cfc8950d87
% canonical-lessons: RU-C317-imeninnik, RU-C317-imeniny, RU-C317-zagadyvat, RU-C317-zaduvat, RU-C317-veselitsya

\chapter{The birthday party}
\label{ch:ru-the-birthday-party}
\hlchaptermodality{317}

\begin{chapteropening}
\textbf{By the end of this chapter:} \emph{I can talk about the birthday boy and a name day, and say that someone made a wish, blew out the candles and had fun.}

Complete the last lesson of chapter 317: I can talk about the birthday boy and a name day, and say that someone made a wish, blew out the candles and had fun.
\end{chapteropening}

\section[\texorpdfstring{\ru{именинник} (imenínnik)}{именинник (imenínnik)}]{\ru{именинник} (imenínnik) — the birthday boy, the person celebrating}

\label{lesson:RU-C317-imeninnik}



\begin{quote}
Before the new one: say the Russian for a procession, then the Russian for a fair.
\end{quote}

\subsection*{You'll want to know: \ru{именинник}}
\textbf{\ru{именинник}} (\emph{imenínnik}) — \textquotedblleft{}the birthday boy, the person celebrating\textquotedblright{}.

\begin{quote}
\ru{Именинник} \ru{получил} \ru{много} \ru{подарков}.
\end{quote}

\begin{quote}
\emph{Imenínnik poluchíl mnógo podárkov.}
\end{quote}

\begin{quote}
The birthday boy got a lot of presents.
\end{quote}

A girl or woman is \emph{\ru{именинница}}.

\begin{cousinweb}
\emph{\ru{именинник}} is built on \emph{\ru{именины}} \textquotedblleft{}a name day\textquotedblright{}, from \emph{\ru{имя}} \textquotedblleft{}name\textquotedblright{}, stem \emph{\ru{имен}-}, which goes back to the PIE word for a name (EtymDB), the source of Latin \emph{nōmen} as well.
\end{cousinweb}

\begin{grammarlens}[title={What you've built}]
One more word for telling, explaining and describing at B1.
\end{grammarlens}

\subsection*{Your turn}
\begin{itemize}
\raggedright
  \item \emph{Say it:} \emph{imenínnik}
  \item \emph{Say it:} \emph{imenínnik}, once more
  \item \emph{Say it:} a sentence of your own with \emph{imenínnik}
  \item \emph{Recall:} say the Russian for a procession, then the Russian for a fair, then say \emph{imenínnik} again
\end{itemize}

\begin{tcolorbox}[breakable,colback=teal!4,colframe=teal!35!black,title={Before you move on}]
What does \ru{именинник} mean? (\textquotedblleft{}The birthday boy, the person celebrating\textquotedblright{}.) Say it once more.
\end{tcolorbox}
\section[\texorpdfstring{\ru{именины} (imeníny)}{именины (imeníny)}]{\ru{именины} (imeníny) — a name day}

\label{lesson:RU-C317-imeniny}



\begin{quote}
Before the new one: say the Russian for a fair, then the Russian for the birthday boy.
\end{quote}

\subsection*{You'll want to know: \ru{именины}}
\textbf{\ru{именины}} (\emph{imeníny}) — \textquotedblleft{}a name day\textquotedblright{}.

\begin{quote}
\ru{Сегодня} \ru{у} \ru{мамы} \ru{именины}.
\end{quote}

\begin{quote}
\emph{Sevódnya u mámy imeníny.}
\end{quote}

\begin{quote}
Today is Mum's name day.
\end{quote}

Plural only: the feast of the saint whose name you bear.

\begin{cousinweb}
\emph{\ru{именины}} is built on the stem \emph{\ru{имен}-} of \emph{\ru{имя}} \textquotedblleft{}name\textquotedblright{} (EtymDB), a cousin of Latin \emph{nōmen} and English \emph{name}.
\end{cousinweb}

\begin{grammarlens}[title={What you've built}]
One more word for telling, explaining and describing at B1.
\end{grammarlens}

\subsection*{Your turn}
\begin{itemize}
\raggedright
  \item \emph{Say it:} \emph{imeníny}
  \item \emph{Say it:} \emph{imeníny}, once more
  \item \emph{Say it:} a sentence of your own with \emph{imeníny}
  \item \emph{Recall:} say the Russian for a fair, then the Russian for the birthday boy, then say \emph{imeníny} again
\end{itemize}

\begin{tcolorbox}[breakable,colback=teal!4,colframe=teal!35!black,title={Before you move on}]
What does \ru{именины} mean? (\textquotedblleft{}A name day\textquotedblright{}.) Say it once more.
\end{tcolorbox}
\section[\texorpdfstring{\ru{загадывать} (zagádyvat)}{загадывать (zagádyvat)}]{\ru{загадывать} (zagádyvat) — to make (a wish); to set (a riddle)}

\label{lesson:RU-C317-zagadyvat}



\begin{quote}
Before the new one: say the Russian for the birthday boy, then the Russian for a name day.
\end{quote}

\subsection*{You'll want to know: \ru{загадывать}}
\textbf{\ru{загадывать}} (\emph{zagádyvat}) — \textquotedblleft{}to make (a wish); to set (a riddle)\textquotedblright{}.

\begin{quote}
\ru{Девочка} \ru{загадала} \ru{желание}.
\end{quote}

\begin{quote}
\emph{Dévochka zagadála zhelániye.}
\end{quote}

\begin{quote}
The girl made a wish.
\end{quote}

Imperfective; perfective \emph{\ru{загадать}} (\emph{zagadát}). \emph{\ru{Загадать} \ru{загадку}} — to ask a riddle.

\begin{cousinweb}
\emph{\ru{загадывать}} is built on \emph{\ru{за}-} and \emph{\ru{гадать}} \textquotedblleft{}to guess, to tell fortunes\textquotedblright{}, inherited from Proto-Slavic \emph{gadati} (EtymDB).
\end{cousinweb}

\begin{grammarlens}[title={What you've built}]
One more word for telling, explaining and describing at B1.
\end{grammarlens}

\subsection*{Your turn}
\begin{itemize}
\raggedright
  \item \emph{Say it:} \emph{zagádyvat}
  \item \emph{Say it:} \emph{zagádyvat}, once more
  \item \emph{Say it:} a sentence of your own with \emph{zagádyvat}
  \item \emph{Recall:} say the Russian for the birthday boy, then the Russian for a name day, then say \emph{zagádyvat} again
\end{itemize}

\begin{tcolorbox}[breakable,colback=teal!4,colframe=teal!35!black,title={Before you move on}]
What does \ru{загадывать} mean? (\textquotedblleft{}To make (a wish); to set (a riddle)\textquotedblright{}.) Say it once more.
\end{tcolorbox}
\section[\texorpdfstring{\ru{задувать} (zaduvát)}{задувать (zaduvát)}]{\ru{задувать} (zaduvát) — to blow out}

\label{lesson:RU-C317-zaduvat}



\begin{quote}
Before the new one: say the Russian for a name day, then the Russian for to make.
\end{quote}

\subsection*{You'll want to know: \ru{задувать}}
\textbf{\ru{задувать}} (\emph{zaduvát}) — \textquotedblleft{}to blow out\textquotedblright{}.

\begin{quote}
\ru{Он} \ru{задул} \ru{все} \ru{свечи} \ru{на} \ru{торте}.
\end{quote}

\begin{quote}
\emph{On zadúl vse svéchi na tórte.}
\end{quote}

\begin{quote}
He blew out all the candles on the cake.
\end{quote}

Imperfective; perfective \emph{\ru{задуть}} (\emph{zadút}), \emph{\ru{задую}}.

\begin{cousinweb}
\emph{\ru{задувать}} is built on \emph{\ru{за}-} and \emph{\ru{дуть}} \textquotedblleft{}to blow\textquotedblright{}, an inherited Slavic verb; EtymDB lists \emph{\ru{дуть}} with no parent.
\end{cousinweb}

\begin{grammarlens}[title={What you've built}]
One more word for telling, explaining and describing at B1.
\end{grammarlens}

\subsection*{Your turn}
\begin{itemize}
\raggedright
  \item \emph{Say it:} \emph{zaduvát}
  \item \emph{Say it:} \emph{zaduvát}, once more
  \item \emph{Say it:} a sentence of your own with \emph{zaduvát}
  \item \emph{Recall:} say the Russian for a name day, then the Russian for to make, then say \emph{zaduvát} again
\end{itemize}

\begin{tcolorbox}[breakable,colback=teal!4,colframe=teal!35!black,title={Before you move on}]
What does \ru{задувать} mean? (\textquotedblleft{}To blow out\textquotedblright{}.) Say it once more.
\end{tcolorbox}
\section[\texorpdfstring{\ru{веселиться} (veselítsya)}{веселиться (veselítsya)}]{\ru{веселиться} (veselítsya) — to have fun, to enjoy oneself}

\label{lesson:RU-C317-veselitsya}



\begin{quote}
Before the new one: say the Russian for to make, then the Russian for to blow out.
\end{quote}

\subsection*{You'll want to know: \ru{веселиться}}
\textbf{\ru{веселиться}} (\emph{veselítsya}) — \textquotedblleft{}to have fun, to enjoy oneself\textquotedblright{}.

\begin{quote}
\ru{Дети} \ru{веселились} \ru{весь} \ru{вечер}.
\end{quote}

\begin{quote}
\emph{Déti veselílis ves vécher.}
\end{quote}

\begin{quote}
The children had fun all evening.
\end{quote}

Imperfective; perfective \emph{\ru{повеселиться}} (\emph{poveselítsya}).

\begin{cousinweb}
\emph{\ru{веселиться}} is built on \emph{\ru{весёлый}} \textquotedblleft{}merry\textquotedblright{}, inherited from Proto-Slavic \emph{veselŭ} (EtymDB), with the reflexive \emph{-\ru{ся}}.
\end{cousinweb}

\begin{grammarlens}[title={What you've built}]
One more word for telling, explaining and describing at B1.
\end{grammarlens}

\subsection*{Your turn}
\begin{itemize}
\raggedright
  \item \emph{Say it:} \emph{veselítsya}
  \item \emph{Say it:} \emph{veselítsya}, once more
  \item \emph{Say it:} a sentence of your own with \emph{veselítsya}
  \item \emph{Recall:} say the Russian for to make, then the Russian for to blow out, then say \emph{veselítsya} again
\end{itemize}

\begin{tcolorbox}[breakable,colback=teal!4,colframe=teal!35!black,title={Before you move on}]
What does \ru{веселиться} mean? (\textquotedblleft{}To have fun, to enjoy oneself\textquotedblright{}.) Say it once more.
\end{tcolorbox}
